\section{Data}
\label{sec:data}

\subsection{Data Sources}

\subsection{Sample Construction}

\subsection{Key Variables}

\subsection{Summary Statistics}

Table \ref{Table-Sum1} presents the summary statistics for the baseline year.

% Table: Summary Statistics
\begin{table}[H]
\normalsize
\centering
\caption{Summary Statistics: Baseline County Characteristics}\label{Table-Sum1}
\setlength{\tabcolsep}{2pt}
\begin{threeparttable}
\begin{tabular}{@{}lccc@{}}
\toprule
\midrule
         & Treatment   & Control    & \textit{p}-value    \\
        &   group  &  group  &             \\
         & (1)  &  (2) &  (3) \\
\midrule
\multicolumn{4}{l}{Panel A: [Category]} \\
\midrule
[Variable 1] \textit{([units])}        & & & \\
[Variable 2] \textit{([units])}        & & & \\
\midrule
\multicolumn{4}{l}{Panel B: [Category]} \\
\midrule
[Variable 3] \textit{([units])}        & & & \\
[Variable 4] \textit{([units])}        & & & \\
\midrule
\multicolumn{4}{l}{Panel C: [Category]} \\
\midrule
[Variable 5]                           & & & \\
[Variable 6]                           & & & \\
\midrule
Number of [Units] & & & \\
\midrule
\bottomrule
\end{tabular}
\begin{tablenotes}[flushleft]
\item \footnotesize \textit{Note:} The table presents [unit]-level mean characteristics for the baseline year. Column (1) reports the mean characteristics for the treatment group, while Column (2) reports those for the control group. Column (3) provides the p-values from t-tests of equal means.
\end{tablenotes}
\end{threeparttable}
\end{table}
